\section{COTS Development Bench}
\label{sec:cots-bench}

Before any custom Alfred board is trusted, build a known-good T1S reference network entirely from commercial hardware. This is Revision~0 -- not an Alfred PCB.

\begin{diagram}
                         Development PC
                              |
                             USB
                              |
                    EVB-LAN8670-USB
                              |
             ==============================
                       10BASE-T1S
             ==============================
                   |                 |
                   v                 v

               LAN8651          LAN8651
               Click board      Click board
                   |                 |
                Clicker        optional node
\end{diagram}

The PC-side EVB-LAN8670-USB adapter gives an independent observation point that does not depend on any Alfred firmware or Alfred hardware -- if a problem appears, this bench isolates whether it is a network/topology issue or an Alfred-specific one.

\subsection{Procure immediately}

\begin{itemize}[itemsep=0.1em]
\item Multiple ($\geq$3) LAN8651-based MikroE Two-Wire ETH Click boards (MIKROE-5543 and/or MIKROE-6550) -- enough to validate a multidrop segment before Alfred custom hardware arrives.
\item Microchip EVB-LAN8670-USB.
\item Required Clicker~4 hardware/accessories.
\item CAN/CAN-FD bench hardware (reference analyzer, bench transceivers, cabling).
\item LAN8660/LAN8661 samples or evaluation hardware, \emph{if} available through Microchip Direct/distribution at time of ordering (Section~\ref{sec:hardware-roles} notes that no dedicated EVB was found for these parts as of the prior research pass -- confirm current availability before assuming none exists).
\end{itemize}

\subsection{Request from Microchip immediately (Week~1 critical-path dependency)}

\begin{itemize}[itemsep=0.1em]
\item LAN8660 secure/NDA datasheet access.
\item LAN8661 secure/NDA datasheet access.
\item Reference schematics and reference layouts for both parts.
\item BOMs for any reference design.
\item Errata.
\item RCP (Remote Control Protocol) documentation.
\item MPLAB Network Creator (tool access/license).
\item Example RCP configurations for control and lighting endpoint use cases.
\item Engineering samples.
\item Endpoint evaluation boards, if any exist beyond what public search identified.
\item Field Applications Engineer (FAE) assistance -- specifically to validate the actuator/LED-driver interface assumptions in Sections~\ref{sec:lan8660-validation}/\ref{sec:lan8661-validation} before Rev~A schematic freeze.
\end{itemize}

\begin{risknote}[title={This is a real, unbounded-duration dependency}]
Everything in Sections~\ref{sec:rev-a-control}--\ref{sec:rev-a-light} depends on this access arriving with enough lead time to inform a schematic. Microchip's response time to an NDA/secure-documentation request is not something this plan can schedule with confidence -- treat it as the single highest-leverage phone call/email of Week~1, escalate immediately if no response by end of Week~1, and have a documented fallback (Section~\ref{sec:risks}) for designing Rev~A from public sell-sheet-level information alone if secure access has not arrived by the Rev~A schematic deadline.
\end{risknote}
